\documentclass{article}
\usepackage[UTF8]{ctex}
\usepackage{amsmath,amssymb,amsfonts}
\usepackage{enumitem}
\usepackage{geometry}
\geometry{a4paper, margin=1in}
% === 错题本排版优化（不参与题目标识） ===
\setlength{\parindent}{0pt}
\setlength{\parskip}{0.35em plus 0.10em minus 0.08em}
\setlength{\emergencystretch}{3em}
\setlist{itemsep=0.25em, topsep=0.40em, parsep=0pt}
\allowdisplaybreaks[2]
% === 错题本排版优化结束 ===
\begin{document}
\section*{福州大学2026年高等代数真题}
\section*{110．福州大学2026年高等代数试卷}
一．填空题。\\

\subsection{题目1}
1．（5 分）设 \(A\) 是 \(n\) 阶正交矩阵且是反对称矩阵，则 \(A^{2}=\) \(\_\_\_\_\)．

\subsection{题目2}
2．（5 分）设 \(\varphi:\binom{a}{b} \mapsto\binom{2 a+b}{a+2 b}\) ，则 \(\varphi\) 在基 \(\binom{1}{0},\binom{1}{1}\) 下的矩阵为 \(\_\_\_\_\)．

\subsection{题目3}
3．（5 分）设 \(\alpha=(1,1,1,1)^{\mathrm{T}} \in \mathbb{R}^{4}\) ，子空间

\[
W=\left\{\left(x_{1}, x_{2}, x_{3}, x_{4}\right)^{\mathrm{T}} \mid x_{1}+x_{2}+x_{3}+x_{4}=0\right\} .
\]

则 \(\alpha\) 到 \(W\) 的距离为 \(\_\_\_\_\)．

\subsection{题目4}
4．（5 分）设 \(A\) 是 \(s \times t\) 矩阵，\(B\) 是 \(t \times s\) 矩阵，若 \(\left(\begin{array}{cc}C & A \\ B & O\end{array}\right)\) 可逆，则 \(A\) 的秩为 \(\_\_\_\_\)．

\subsection{题目5}
5．（5 分）设 \(f(x)=\prod_{j=1}^{2026}\left(x^{2}+j\right)^{j^{2}+j+1}, d(x)\) 是 \(f(x)\) 与 \(f^{\prime}(x)\) 的首一最大公因式，则 \(f(x)\) 除以 \(d(x)\)的商为 \(\_\_\_\_\)．

二．简答题。\\

\subsection{题目6}
6．（5 分）计算 \(n\) 阶行列式

\[
\left|\begin{array}{cccc}
1-a_{1}^{2} & -a_{1} a_{2} & \cdots & -a_{1} a_{n} \\
-a_{2} a_{1} & 1-a_{2}^{2} & \cdots & -a_{2} a_{n} \\
\vdots & \vdots & & \vdots \\
-a_{n} a_{1} & -a_{n} a_{2} & \cdots & 1-a_{n}^{2}
\end{array}\right| .
\]

\subsection{题目7}
7．（5 分）设 \(n\) 阶方阵 \(A\) 有 \(n\) 个不同的特征值 \(\lambda_{1}, \lambda_{2}, \cdots, \lambda_{n}\) ，求 \(A\) 的极小多项式。

\subsection{题目8}
8．（5 分）设 \(n\) 阶方阵

\[
A=\left(\begin{array}{ccccc}
1 & -1 & 1 & \cdots & (-1)^{n} \\
& 1 & -1 & \ddots & \vdots \\
& & 1 & \ddots & 1 \\
& & \ddots & -1 \\
& & & 1
\end{array}\right) .
\]

求 \(A\) 的 Jordan 标准型．

\subsection{题目9}
9．（5 分）设 \(A, B\) 是 \(n\) 阶正定矩阵，试问：\(A B\) 是否必为正定矩阵？若是，请说明理由；若否，请给出反例．

\subsection{题目10}
10．（5 分）设 \(A, B\) 是 \(n\) 阶实方阵，且存在 \(n\) 阶复可逆矩阵 \(P\) 使得 \(B=P^{-1} A P\) ，试问：是否存在 \(n\)阶实可逆矩阵 \(Q\) 使得 \(B=Q^{-1} A Q\) ？若是，请说明理由；若否，请给出反例。
\end{document}
